% Part: model-theory
% Chapter: models-of-arithmetic
% Section: introduction

\documentclass[../../../include/open-logic-section]{subfiles}

\begin{document}

\olfileid{mod}{mar}{int}
\section{પરિચય}

અંકગણિતનો \emph{પ્રમાણભૂત નિદર્શ} એ $\Domain{N}=\Nat$ ધરાવતી
!!{structure}~$\Struct{N}$ છે, જેમાં $\Obj{0}$, $\prime$, $+$,
$\times$ અને~$<$નું આપણે અપેક્ષા રાખીએ એ રીતે અર્થઘટન થાય છે. એટલે
$\Obj{0}$ એ~$0$ છે, $\prime$ ઉત્તરગામી વિધેય છે, અને $+$ તથા
$\times$નું અનુક્રમે $\Nat$ની સંખ્યાઓના સરવાળા અને ગુણાકાર તરીકે
અર્થઘટન થાય છે. વધુ ચોક્કસ રીતે,
\begin{align*}
  \Assign{\Obj{0}}{N} & = 0\\
  \Assign{\prime}{N}(n) & = n + 1\\
  \Assign{+}{N}(n, m) & = n + m\\
  \Assign{\times}{N}(n, m) & = nm
\end{align*}
અલબત્ત, $\Lang{L_A}$ માટે એવી સંરચનાઓ પણ છે જેમનાં વ્યાપકક્ષેત્ર
$\Nat$થી જુદાં છે. દાખલા તરીકે, $\Domain{M}=\{a\}^*$ ધરાવતી
$\Struct{M}$ લઈ શકીએ છીએ ($a$ એ એકમાત્ર સંકેત હોય તેવી સાન્ત
શ્રેણીઓ, એટલે કે $\emptyset$, $a$, $aa$, $aaa$, \dots) અને તેનાં
અર્થઘટનો આ પ્રમાણે આપી શકીએ છીએ:
\begin{align*}
  \Assign{\Obj{0}}{M} & = \emptyset\\
  \Assign{\prime}{M}(s) & = s \concat a\\
  \Assign{+}{M}(a^n, a^m) & = a^{n + m}\\
  \Assign{\times}{M}(a^n, a^m) & = a^{nm}\\
  \Assign{<}{M} & = \Setabs{\tuple{a^n,a^m}}{n,m\in\Nat\text{ અને }n<m}
\end{align*}
આ બે સંરચનાઓ ``મૂળતઃ એકસરખી'' છે: તેમનાં !!{domain}sના
!!{element}s જુદા છે, પરંતુ અર્થઘટન-વિધેયો એ !!{domain}sના
!!{element}sને એકબીજા સાથે જે
રીતે જોડે છે તે જુદી નથી. આપણે કહીએ છીએ કે આ બે !!{structure}s
\emph{એકરૂપ} છે.
\sourcecorrection{OLMAR-001}{સ્થિર અંગ્રેજી મૂળની
  \texttt{introduction.tex}, પંક્તિઓ~32--33માં $\{a\}^*$ના ઘટકોને બદલે
  પ્રાકૃતિક સંખ્યાઓ $n,m$ને સરવાળા અને ગુણાકારનાં દલીલખંડો બનાવવામાં
  આવ્યા છે. અહીં ક્ષેત્રના વાસ્તવિક ઘટકો $a^n,a^m$ પુનઃસ્થાપિત કર્યા
  છે.}
\sourcecorrection{OLMAR-002}{એ જ મૂળની પંક્તિઓ~29--34માં
  $\Lang{L_A}$ના સંકેત~$<$નું અર્થઘટન આપવામાં આવ્યું નથી; તેથી આપેલો
  આંકડો પૂર્ણ $\Lang{L_A}$-સંરચના બનતો નથી અને દાવો કરેલી એકરૂપતા પણ
  નક્કી થતી નથી. અહીં $a^n<a^m$ ત્યારે અને માત્ર ત્યારે જ્યારે $n<m$
  એવું અર્થઘટન ઉમેર્યું છે.}

સઘનતા પ્રમેયનું એક સરળ પરિણામ એ છે કે $\Struct{N}$માં સાચો દરેક
સિદ્ધાંત એવા નિદર્શો પણ ધરાવે છે જે $\Struct{N}$ને એકરૂપ નથી. આવી
સંરચનાઓને \emph{અપ્રમાણભૂત} કહે છે. તેમાં રસપ્રદ વાત આ છે: પ્રમાણભૂત
નિદર્શના !!{element}s (એટલે $\Struct{N}$ના, તેમજ તેને એકરૂપ બધી
!!{structure}sના ઘટકો) પ્રમાણભૂત આંકિક પદો~$\num{n}$નાં મૂલ્યોથી
સંપૂર્ણપણે મળી જાય છે, એટલે કે
\[
\Domain{N} = \Setabs{\Value{\num{n}}{N}}{n \in \Nat}
\]
પરંતુ અપ્રમાણભૂત નિદર્શોમાં એવું નથી: જો $\Struct{M}$ અપ્રમાણભૂત હોય,
તો ઓછામાં ઓછો એક એવો $x\in\Domain{M}$ છે કે દરેક~$n$ માટે
$x\neq\Value{\num{n}}{M}$.

આ અપ્રમાણભૂત ઘટકો ખૂબ રસપ્રદ છે: તેઓ ``અનંત પ્રાકૃતિક સંખ્યાઓ'' છે.
પરંતુ એક અર્થમાં તેમનું અસ્તિત્વ અપૂર્ણતાની ઘટનાઓને પણ સમજાવે છે.
દાખલા તરીકે, પેયાનો અંકગણિત માટેનું સુસંગતતા-વિધાન
$\OCon[\Th{PA}]$, એટલે કે
$\lnot\lexists[x][\OPrf[\Th{PA}](x,\gn{\lfalse})]$, વિચારો.
$\Th{PA}$ ન તો $\OCon[\Th{PA}]$ને સાબિત કરે છે, ન
$\lnot\OCon[\Th{PA}]$ને; તેથી
બેમાંથી કોઈને પણ $\Th{PA}$માં સુસંગત રીતે ઉમેરી શકાય છે. $\Th{PA}$
સુસંગત હોવાથી $\Sat{N}{\OCon[\Th{PA}]}$, અને પરિણામે
$\Sat/{N}{\lnot\OCon[\Th{PA}]}$. તેથી $\Struct{N}$ એ
$\Th{PA}\cup\{\lnot\OCon[\Th{PA}]\}$નો નિદર્શ \emph{નથી}, અને તેના
બધા નિદર્શો અપ્રમાણભૂત હોવા જોઈએ.
$\Th{PA}\cup\{\lnot\OCon[\Th{PA}]\}$ના નિદર્શમાં એવું કોઈ
!!{element} હોવું જોઈએ જે
$\lexists[x][\OPrf[\Th{PA}](x,\gn{\lfalse})]$ને સત્ય બનાવતા સાક્ષીનું
કામ કરે—એટલે $\Th{PA}$માંથી વિરોધાભાસની !!a{derivation}ની ગેડલ સંખ્યા.
આવો !!a{element} પ્રમાણભૂત હોઈ શકે નહિ, કારણ કે દરેક~$n$ માટે
$\Th{PA}\Proves\lnot\OPrf[\Th{PA}](\num{n},\gn{\lfalse})$.
\sourcecorrection{OLMAR-003}{સ્થિર અંગ્રેજી મૂળની
  \texttt{introduction.tex}, પંક્તિ~67માં અસ્તિત્વવાચક રીતે બંધાયેલો
  ચલ~$x$, દ્વિસ્થાની સાબિતી-વિધેય $\OPrf[\Th{PA}]$ના દલીલખંડમાંથી
  ગાયબ છે. પંક્તિ~59ના સુસંગતતા-સૂત્ર સાથે સુસંગત રીતે અહીં
  $\OPrf[\Th{PA}](x,\gn{\lfalse})$ પુનઃસ્થાપિત કર્યું છે.}

\end{document}
